\begin{table}[H]
\centering
\caption{Formas de $y_{h}$ en base a los tipos de raiz}
\begin{threeparttable}
{\setstretch{1.75}
\begin{tabular}{p{0.45\linewidth}p{0.45\linewidth}}\hline
    
    \textbf{Tipo de raiz} & \textbf{Forma de $y_{h}$} \\ \hline
    
    Reales y distintas ($\lambda_{1} \neq \lambda_{2}$) & $C_{1}e^{\lambda_{1}x} + C_{2}e^{\lambda_{2}x}$ \\
    Reales e iguales ($\lambda_{1} = \lambda_{2}$) & $(C_{1} + C_{2} \cdot x) \cdot e^{\lambda x}$ \\
    Complejas ($\lambda = \alpha \pm \beta i$) & $e^{\alpha x} \cdot (C_{1}\cos{(\beta x)} + C_{2}\sin{(\beta x)})$ \\
    
    \hline
\end{tabular}
}
\begin{tablenotes} % ex: \tntote{1} ... \item[1]
    \item[] 
\end{tablenotes}
\end{threeparttable}
\label{tab:eq_segundo_orden}
\end{table}